\documentclass[11pt,a4paper]{article}
\usepackage[utf8]{inputenc}\usepackage[T1]{fontenc}\usepackage{lmodern}
\usepackage{amsmath,amssymb,amsthm}\usepackage{geometry}\usepackage{xcolor}
\usepackage{hyperref}\usepackage{microtype}\usepackage{enumitem}\usepackage{booktabs}
\geometry{margin=2.5cm}
\hypersetup{colorlinks=true,linkcolor=blue!55!black,urlcolor=blue!55!black}
\theoremstyle{plain}
\newtheorem{theorem}{Theorem}[section]
\newtheorem{lemma}[theorem]{Lemma}
\newtheorem{corollary}[theorem]{Corollary}
\theoremstyle{definition}
\newtheorem{remark}[theorem]{Remark}
\sloppy

\title{\bfseries Closed-Form Orbits of the Maps $ax+1$:\\ A Characterisation of Fermat Bases}
\author{Joel Cruz Cabrera\\
\small Santa Cruz de Tenerife, Canary Islands, Spain\\
\small\texttt{tf.joel.work@gmail.com}\quad ORCID: 0009-0005-4048-1237}
\date{July 2026}

\begin{document}\maketitle

\begin{abstract}
For odd $a$ consider the accelerated map $x\mapsto\operatorname{odd\_part}(ax+1)$ on odd
integers, and the associated ``block'' parametrisation $x=2^{r+1}m-1$. We characterise
exactly those bases $a$ for which the iterates admit the closed polynomial form
$x_i=(a^i m-1)/(a-1)$: this happens \emph{if and only if} $a-1$ is a power of $2$, i.e.
$a=2^e+1$ --- a Fermat-type base ($3,5,9,17,33,\dots$). Collatz is the case $a=3$. Both
directions follow from a single one-step law, valid for every odd $a$. \textbf{The whole
characterisation --- the one-step law, the biconditional, and the orbit formula --- is
formalised in Lean~4 and machine-checked with kernel-only axioms.}
\textbf{This is a theorem about iteration of affine maps and is independent of the Collatz
conjecture.}
\end{abstract}

\noindent\textbf{MSC:} 37P05 (Arithmetic dynamics), 11B37 (Recurrences), 68V20
(Formalisation of mathematics).

\section{Setting}
Fix an odd integer $a\ge3$. On the odd integers define
\[
  \operatorname{step}_a(x)=\frac{ax+1}{2^{v_2(ax+1)}},
\]
where $v_2$ is the $2$-adic valuation. Every odd $x$ can be written uniquely as
\begin{equation}\label{eq:block}
  x=2^{r+1}m-1,\qquad m \text{ odd},\ r\ge0,
\end{equation}
so that $r=v_2(x+1)-1$ records how many leading valuation-one steps the orbit begins with.

The question we answer is: \emph{for which $a$ do the iterates of $\operatorname{step}_a$
starting from \eqref{eq:block} follow a single closed-form polynomial law?}

\section{The characterisation}

\begin{theorem}[One-step law, and the characterisation]\label{thm:main}
Let $a\ge3$ be odd and write $a-1=2^{e}q$ with $q$ odd. For every odd $m$ and every $r\ge e$,
put $x=2^{r+1}m-1$. Then
\begin{equation}\label{eq:onestep}
  v_2(ax+1)=e\qquad\text{and}\qquad \operatorname{step}_a(x)=a\,2^{\,r+1-e}m-q .
\end{equation}
Hence the closed-form law $\operatorname{step}_a(x)=a\,2^{\,r+1-e}m-1$ holds \emph{if and
only if} $q=1$, i.e.\ iff $a-1$ is a power of $2$, i.e.\ iff $a=2^{e}+1$. In that case
iteration gives
\begin{equation}\label{eq:cf}
  x_i=a^{i}\,2^{\,r+1-ie}m-1\qquad\text{for every }i\text{ with }ie\le r .
\end{equation}
\end{theorem}

\begin{proof}
Compute:
\[
  ax+1=a(2^{r+1}m-1)+1=a\,2^{r+1}m-(a-1)=2^{e}\bigl(a\,2^{\,r+1-e}m-q\bigr).
\]
As $r\ge e$, the bracket is (even)$-$(odd), hence odd; so $v_2(ax+1)=e$ \emph{exactly} and
$\operatorname{step}_a(x)$ is the bracket, which is \eqref{eq:onestep}. Its value equals
$a2^{\,r+1-e}m-1$ precisely when $q=1$ --- both directions at once.

For \eqref{eq:cf}, let $q=1$. The image $a2^{\,r+1-e}m-1$ has again the shape
\eqref{eq:block}, with odd parameter $m'=am$ and exponent $r'+1=r+1-e$; so the family is
preserved and induction on $i$ gives \eqref{eq:cf} while $ie\le r$.
\end{proof}

\noindent The one-step law \eqref{eq:onestep} holds for \emph{every} odd $a$; the
characterisation is simply the observation that its right-hand side falls back into the
family exactly when the odd part $q$ of $a-1$ is trivial.

\begin{table}[ht]\centering
\begin{tabular}{c c c c}
\toprule
$a$ & $a-1=2^{e}q$ & $q$ & closed form?\\\midrule
$3$  & $2^1\cdot1$ & $1$ & yes (Collatz)\\
$5$  & $2^2\cdot1$ & $1$ & yes\\
$7$  & $2^1\cdot3$ & $3$ & no\\
$9$  & $2^3\cdot1$ & $1$ & yes\\
$11$ & $2^1\cdot5$ & $5$ & no\\
$13$ & $2^2\cdot3$ & $3$ & no\\
$15$ & $2^1\cdot7$ & $7$ & no\\
$17$ & $2^4\cdot1$ & $1$ & yes\\
$33$ & $2^5\cdot1$ & $1$ & yes\\
\bottomrule
\end{tabular}
\caption{The closed form holds exactly when the odd part $q$ of $a-1$ is $1$. The one-step
law and the closed form were tested independently by running the map --- the ``no'' column is
a genuine test, not a restatement of the hypothesis --- for all odd $a\le999$.}
\end{table}

\begin{remark}[Terminology]
The bases $a=2^e+1$ are $3,5,9,17,33,65,129,257,513,\dots$; we call them \emph{Fermat-type}
because the Fermat numbers $2^{2^k}+1$ are the sub-family $e=2^k$. The two sets are
\emph{not} equal: $9,33,65,129,513$ are of the form $2^e+1$ but are not Fermat numbers.
\end{remark}

\section{Formalisation}
The orbit formula \eqref{eq:cf} is formalised in Lean~4 as \texttt{orbit\_formula\_fermat}:
\[
  \texttt{iter\_a}\ a\ i\ (2^{r+1}m-1)=a^{i}\,2^{\,r+1-ie}m-1\qquad(ie\le r),
\]
for $a=2^{e}+1$, $e\ge1$, $m$ odd. It builds with
\begin{center}
  \texttt{lake build CollatzProof.KodamaFermat}
\end{center}
and \texttt{\#print axioms} reports exactly the three ambient kernel axioms
(\texttt{propext}, \texttt{Classical.choice}, \texttt{Quot.sound}) --- no \texttt{sorry} and
no custom axioms. The audit lives in \texttt{CollatzProof/AxiomAudit.lean} so that it re-runs
on every build.

Since 2026-07-17 the \emph{whole} characterisation is in the kernel, not just one half:
\begin{itemize}[nosep]
  \item \texttt{step\_a\_general}: the one-step law \eqref{eq:onestep} for \emph{every} odd
        base $a$ --- write $a-1=2^{e}q$ with $q$ odd, then $v_2(an+1)=e$ and
        $\operatorname{step}_a(2^{s+e}m-1)=a2^{s}m-q$;
  \item \texttt{closed\_form\_iff\_fermat}: the biconditional
        $\operatorname{step}_a(2^{s+e}m-1)=a2^{s}m-1 \iff q=1$;
  \item \texttt{fermat\_iff\_pow\_two}: $a=2^{e}+1 \iff q=1$.
\end{itemize}
All report exactly the three ambient kernel axioms. The audit
\texttt{CollatzProof/AxiomAudit.lean} re-runs \texttt{\#print axioms} on every one of them at
each build, so the claim cannot silently rot.

\section{Consequences and scope}
\begin{corollary}
The Collatz map ($a=3$) is the smallest Fermat-type base; the next are $5,9,17,33,\dots$ (all
such bases are odd, $2^e+1$ being odd for $e\ge1$).
The heuristic drift of $\operatorname{step}_a$ is $\lambda_a=\log a-\mathbb{E}[v_2]\log2
=\log a-2\log2$, using $\mathbb{E}[v_2]=2$ from the companion note~\cite{geom},
negative iff $a<4$; thus $a=3$ is the unique odd base on the contracting side
and every odd $a\ge5$ is super-critical.
\end{corollary}

\begin{remark}[What this does \emph{not} say]
Theorem~\ref{thm:main} is a statement about the algebraic shape of the iterates. It does not
decide whether any particular orbit converges or diverges: the drift $\lambda_a$ is an
average and, as the divergence example $\lfloor\sqrt2\,n+4\rfloor$ of
Alth\"ofer--Zipproth shows, the sign of an average drift does not settle the fate of a
deterministic orbit. In particular this paper says nothing about the Collatz conjecture, nor
about the open question of whether $9x+1$ diverges from $x=1$.
\end{remark}

\section*{Reproducibility}
The base table is regenerated by \texttt{tools/verify\_for\_papers.py}; the Lean development
builds with \texttt{lake build}.

\section*{Acknowledgement}
AI-assisted; every claim was verified by exact computation or Lean's kernel before inclusion.

\begin{thebibliography}{9}
\bibitem{geom} J.~Cruz Cabrera, \emph{An exact equidistribution law for the $2$-adic
valuations of $3^rm-1$}, companion note, 2026.
\bibitem{lagarias} J.~C. Lagarias, \emph{The $3x+1$ problem and its generalizations},
Amer. Math. Monthly \textbf{92} (1985), 3--23.
\bibitem{mathlib} The mathlib Community, \emph{The Lean mathematical library}, CPP 2020.
\bibitem{az} I.~Alth\"ofer, T.~Zipproth, \emph{Divergence in a Collatz-type recursion},
preprint, August 2025.
\end{thebibliography}
\end{document}
